\section{Introduction}

Locality in lattice gauge theory is deceptively representation-dependent.  A
plaquette action is local in link variables, while the part of that action
seen by one link is naturally expressed through the surrounding staples.  In
the staple description, changing the central link looks like a one-point
optimization.  In the link description, however, changing the relative
positions of the staples requires modifying shared links and therefore
changes plaquettes outside the chosen star.  This paper makes that mismatch
quantitative.

Staples are central to heat-bath and overrelaxation updates, cooling, link
smearing, and gauge-field flow.  These methods exploit the dependence of a
local action on a link and its surrounding paths; see, for example,
\cite{Adler1981,BrownWoch1987,MorningstarPeardon2004,Luscher2010}.  Our aim is
not to propose another smoothing algorithm.  Instead, we isolate an invariant
quantity that records how mutually compatible the surrounding transports are,
and then ask whether an ideal operation on that quantity is realizable by an
actual gauge field.

The first observation is nonlinear and coordinate-free.  If $S_1,\ldots,S_m$
are the complementary transports around a link, their Fr\'echet variance
\[
  D(S_1,\ldots,S_m)=\inf_C\sum_{i=1}^m d(C,S_i)^2
\]
is gauge invariant even though neither $S_i$ nor $C$ is.  It is precisely the
minimum local geodesic-quadratic action obtainable by varying the central
link.  On $\SU\simeq S^3$, a comparison with tangent-space variance then
relates $D$ to the average of the pairwise squared distances.

The second observation concerns dependence.  In the linearized Gaussian
theory, plaquette curvatures are not independent Gaussians: their covariance
is the orthogonal projection onto the image of the link-to-plaquette
coboundary.  Restricting and centering this projection produces a nontrivial
five-mode law.  We evaluate it exactly on the smallest open four-dimensional
block with interior links.

The third observation reverses the map.  Suppose one prescribes a centered
change of the six curvatures incident on a link.  Among all link variations
that realize it, minimize the curvature created on the other plaquettes.  The
answer is a positive quadratic form.  Its smallest eigenvalue measures the
irreducible spillover of a staple-space operation when lifted back to link
space.  On our block it is $37/33>1$: there is no nonzero external-flat
repair mode.

These three results fit into one diagram:
\begin{center}
\begin{tikzpicture}[node distance=16mm and 23mm,>=Latex,every node/.style={align=center}]
\node (links) [draw,rounded corners] {link variations\\$C^1$};
\node (curv) [draw,rounded corners,right=of links] {plaquette curvature\\$d_1 C^1\subset C^2$};
\node (star) [draw,rounded corners,above right=8mm and 22mm of curv] {centered star\\dispersion};
\node (ext) [draw,rounded corners,below right=8mm and 22mm of curv] {external\\curvature};
\draw[->] (links)--node[above] {$d_1$} (curv);
\draw[->] (curv)--node[above left] {$C_m O_\ell S_\ell$} (star);
\draw[->] (curv)--node[below left] {$E_\ell$} (ext);
\draw[<-,dashed] (links) to[bend left=16] node[below] {minimum-spill lift} (star);
\end{tikzpicture}
\end{center}
The forward arrows expose an observable; the dashed reverse arrow is a
constrained reconstruction problem.  The spillover form is the obstruction
to treating that reconstruction as lossless.

The paper is organized as follows.  Sections~\ref{sec:setup}--\ref{sec:nonlinear}
develop the gauge-covariant nonlinear dispersion.  Section~\ref{sec:gaussian}
derives its exact Gaussian law on the finite block.  Section~\ref{sec:spill}
defines the general minimum-spill form and evaluates it.  We then record the
consequences for repair budgets and for representation change.  The
appendices contain complete proofs and the exact certificate contract.
